\begin{tikzpicture}[
  anak/.style={draw,rounded corners,align=center,text width=1.68cm,
    minimum height=1.05cm,font=\scriptsize,fill=blue!8},
  panah/.style={-{Stealth[length=2mm]},thick}
]
  \node[anak,fill=gray!35] (t0) at (0,0)
    {\textbf{0}\\{\tiny\mdseries belum memadai}};
  \node[anak] (t1) at (1.95,0.8)
    {\textbf{1}\\{\tiny\mdseries aktivitas belum lengkap}};
  \node[anak] (t2) at (3.9,1.6)
    {\textbf{2}\\{\tiny\mdseries aktivitas dasar lengkap}};
  \node[anak,fill=orange!12] (t3) at (5.85,2.4)
    {\textbf{3}\\{\tiny\mdseries memakai aset organisasi}};
  \node[anak] (t4) at (7.8,3.2)
    {\textbf{4}\\{\tiny\mdseries kinerja diukur kuantitatif}};
  \node[anak,fill=green!10] (t5) at (9.75,4.0)
    {\textbf{5}\\{\tiny\mdseries perbaikan berkelanjutan}};
  \draw[panah] (t0.north) |- (t1.west);
  \draw[panah] (t1.north) |- (t2.west);
  \draw[panah] (t2.north) |- (t3.west);
  \draw[panah] (t3.north) |- (t4.west);
  \draw[panah] (t4.north) |- (t5.west);
  \node[font=\scriptsize,align=left,text width=6.6cm] at (6.75,0.3)
    {Tiap tingkat diperiksa lewat aktivitas dan buktinya pada cakupan yang tertulis.
    Tingkat yang lebih tinggi tidak otomatis menjadi target semua objektif.};
\end{tikzpicture}
